\section{声明纪律与负结果报告规范}
\label{sec:discipline}

实验扩容最大的风险不是跑不完，而是``数据多了之后挑对自己有利的说''——这恰是实验分析规范（awesome-ai-research-writing 仓库的实验分析 prompt 所内嵌的原则：所有结论必须严格基于输入的数据；严禁编造数据、夸大提升幅度；如果数据中没有明显的优势或趋势，如实描述，不要强行总结显著提升）所禁止的。本节把声明纪律前置为实验设计的一部分，规则如下。

第一，先声明后实验。每个实验单元在运行前登记其``将支撑的论文声明''与``判定标准''（本计划各节已给出，执行时可对照），实验完成后无论结果方向如何，登记的声明按原判定标准兑现或撤回。例如 A1 的登记是：若 \sieve{}--head 配对区间半宽收窄至 $\pm 5$ 点内且区间跨零，则论文声明维持``统计持平''并加注 $n$；若区间分离，声明升级为方向性优势——两种结局都算完成，但措辞跟着证据走。

第二，负结果是发现，不是失败。BM25 召回更高的既有叙事已经示范了这一点；扩容后的负结果同样按``结构性发现''报告：B1 若出现阅读器异号，报告为``阅读器家族与压缩敏感性的交互''；A2 若多跳数据集端到端退化快于召回，报告为``召回--质量传导率的数据集依赖''，并与相关研究节互相印证。唯一禁止的写法是把负结果藏进附录或不报。

第三，协议冻结与版本登记。提示模板、解码参数、预算定义、池种子在每条路径内冻结；任何必须的改动（例如阅读器换代的 API 不再可用）单独登记为协议版本号，新旧数据不跨版本合并统计。端到端协议的冻结条款已在路径 A 声明，B/C 同理。

第四，宏引用纪律延续。第一阶段审计发生过一次凭印象写错消融数字（43.6 误写为 44.6）的事故并被 numbers.tex 宏核对纠正；扩容后数字数量翻倍，一切进正文的数字必须经 \texttt{make\_numbers\_journal.py} 生成的宏，reviewer 复核脚本同时校验宏值与结果 JSON 的一致性。

第五，与 anonymity 及投稿伦理的衔接。扩容实验全部基于公开基准与自建池，无新数据许可问题；多阅读器 API 使用遵守各服务商条款并如实披露阅读器版本与解码参数；本地阅读器权重与推理配置随仓库发布，保证可复现声明不因阅读器多样性而弱化。
